\daymark{8}{Stone 14, Y.H. 812 (Fourthday)}

Nothing in the house had told him they were going. \person{Sella Vail} took the basket down from its peg in the half-dark while the coals still lay under their ash, and the little-door burn closed first on the handle and then on the fold, and she said to the wool what she said to the pot: “Keep it off his ears. It's the river wind.” Eight days old, carried because there was nobody in the yard to leave him with, he went out through the door that did not quite shut and down the \place{Clay Yard} with the cold coming up off the \place{the Slip Brook} to meet him.

She carried him the way she carried a pot filled that must not be slopped, her forearm under the basket's belly, the burn shut around the handle. His head went where it wanted, toward the strip of sky between the roofs, and when it went too far she put it back with two fingers behind the ear without once looking at his face, and that was the whole of her attention and the whole of her refusal in one movement. The yard was trodden ash and frost, and the frost lay over the ash in a skin that broke under her heel. The drying racks stood with their empty slats up. The wheel-shed was dark behind him. The house went on breathing four, and four, and the breath it missed, long after the wall was between them, and he kept the count going after there was nothing to count, because counting was a thing a man could do with his hands full and there was no part of him that had ever been able to stop.

The walk took the time a walk always took and Sella did not hurry it. The Slip Brook went over its stones beside the path with the thin skin on it that the children of \place{Low Kiln} read as weather, and a woman up to her elbows in a tub at a door looked up as the basket passed and went back to her work without appearing to have seen it, which was the whole attention the district had to give a woman carrying a bundle. Sella kept to the middle of the path and a half-pace from every door where a slate might lie open, and the boy, who had learned the half-pace in the market at \place{the Cross} and had never once been told what it was for, counted the doors instead. He counted seven and lost the count at the eighth when his head went where it wanted and the small sound came out of him and Sella put two fingers behind the ear and set him right. Seven, and a place where the eighth should have been. It was becoming a habit of his, this finding a gap where a thing should have been, and he had no standing to call it anything.

The town woke in layers on the walk down. First the gulls, crying above the roofs before there was light enough to see them, their cry coming down thin and peeled; then the smell of wet rope and salted perch; then the small sound of water moving under ice somewhere ahead; then the cold, which came last and stayed. By the time the steps of the \place{the Fishstairs} were under Sella's feet the whole place was made of one material, and that material was cold, and the cold was a substance a body could put a hand against. He knew that much without being told it. He had spent eight days learning that warmth in this country was always somebody's, entering and leaving rooms the way a person enters and leaves a room, and that the cold was the part nobody entered and the part that kept whatever was taken. He lay with his cheek against the fold and felt it come in at ear and temple and wrist, in that order, which was an order he had learned to know without ever having chosen to learn it.

At the foot of the steps the market had not opened in earnest but was already full of the things that would be sold later: a man setting out a board of perch with his back to the wind, two girls turning a barrow of wet weed, a collector's boy with a satchel waiting for nothing in particular. Rumours went around before the fish did, though they would not be priced and cried until noon; the boy heard a figure and then another and could not spend either of them, and the language they were spoken in was a row of sounds he had absorbed through his ears and had never been given a key to. It was an account read aloud in a tongue he did not have, and he was carried through it at the level of Sella's hip with his mouth working under the wool. A brazier of coals stood by the eel stall with a coal at its edge that had fallen out and gone black on the stone, and the boy, who knew what a banked fire was worth in this district, noted that even here warmth was kept and counted and lost by inches.

The \place{the north bridge} stood over them with its shadow on the steps. The \place{the Ilm} lay under its skin of ice, a skin the colour of skimmed milk left standing too long, and the ice held, and the district said it held, and nobody on the steps looked at it twice, because there was nothing a body could do with ice. Downstream, beyond the bridge, the river went over the ledge of hard stone that had given the town its name, the \place{the Singing Sill}, and gave back a low note when \lore{the Warm} rose in it; he had the name by listening and had never had it explained, and he set it down in the file he kept among the things that were named and not accounted for. A \lore{brine-gull} came down over the landing with something bright in its beak, a coal not yet cold from somebody's brazier, and went up again with the gulls crying after it, and the boy noted that warmth here was a thing that could be stolen in flight and carried off and eaten, which was not a thought he enjoyed and was nevertheless exact.

At the bottom of the steps the stall was already open and \person{Mira Silt} was in it up to the elbows. Slime had dried pale on her sleeves in stiff folds the colour of old bone, and she had the salt of the perch up to her wrist and river silt in the creases of her neck, and she smelled of silt and tallow smoke and the cold off the water all at once, which was how he had learned to smell her in two rooms before he had ever seen her behind a stall. The top of her left ear was gone. She held her head a little to the right against the cold in the place where it had been, the way a man holds his head away from a draught, and she said a hook had taken it, and the scar said a knife. He read the scar the way he had once read a wound on a man who told him he had fallen: a clean edge drawn down and across, not the ragged tear a hook would leave. The small smile sat on her face the way a door sits on the latch, able to be shut without anybody seeing the movement. She kept it open for the stall and for whoever came down the steps.

She had three \lore{eel-bones} in her left hand and she flicked them against her thumb while she looked the morning over, one, and another, and the third, deciding something, and when the deciding was done she put one away in her pocket and went on with the other two. He had watched her do that in the birth room and at the edge of the Cross and had never once seen her put the bone away before the deciding was finished, so that the movement was a tell, and he had nobody to tell it to. She worked with a short knife that had been sharpened on a stone and not cleaned, and the eels came up out of a box of wet weed and went down under her hand in two passes, and the bones went into a bowl at her right and came out again when she wanted to think. Everything about her was quick and had been quick for fifteen years. Nothing about her was careless.

The knife went down the length of an eel and came up clean where the hand had been, and the bones came away in the bowl with a small sound like rain on a board, and the boy, who had once stood at a table and separated the parts of a body with an instrument he had chosen himself, recognised the economy of the movement: nothing wasted, nothing handled twice. Her hands were broad and worked and the nails cut short to the quick, and she never once cut herself, and she never once looked into the basket at her elbow. That last was the most difficult movement of all and she made it look like nothing. She talked to the eels as she worked and to the woman at the perch tub and to nobody in the basket, and the boy lay in the basket and understood perfectly that he was being kept by a person who would not acknowledge him in front of a slate. He had no word for the understanding and no use for it, and he kept it.

\scenebreak

\person{Berr Clay} came down the steps with the red wool cord of his office folded away somewhere under his coat and his hands large and clean. He wiped his thumb on his apron before he touched the slate he carried, so that the slate would stay clean, and the ring of tasted glaze showed green at one corner of his mouth. He had the smell of hot clay on him and, under it, the smell of nothing at all, which the boy had learned to recognise as the most expensive smell in \place{Grath}. The green at the mouth was the colour of the glaze the potters tasted to judge it, and the boy knew a man in the old life who had taken something sweet into him for thirty years and had a pale line at his gum to show for it, and he set the two men side by side and joined neither of them.

“\lore[the levy]{The levy} is called for the feast,” Berr Clay said to the eels. “It is not called by me.”

“Two marks the pair or two marks the pair and two short?” Mira said. “Don't trouble yourself. It's two marks the pair and two short, and I've counted them twice, and I'll count them again after I've eaten.”

He did not answer, because she had left nothing in the sentence that wanted an answer. He wiped his thumb once more on his apron, which was already clean, and set two marks on the corner of his slate, and she put four eels into his basket where six should have gone, and both of them went on with the morning as if nothing at all had passed between them. Six years of it, and neither had ever said a word about the two eels or the two marks. The boy had once watched a nurse and a billing clerk settle a question about a consumable with the same careful avoidance of the subject, two people standing close and speaking around a thing in front of a man who could not ask, and he took the shape of it in now without a word and without letting anyone see he had taken it. Berr went back up the steps with the slate held flat and the green at his mouth, reasonable to the last, and the gulls went over him.

Sella had not set the basket down. She stood a half-pace from the stall board where Berr's slate had lain open, and when the slate was shifted an inch toward her she shifted an inch away, and none of it was visible as a movement anybody could have called a choice. She said to the wool, “Keep it off his ears,” and to the river wind, “You'll pay for that come thaw,” and never once said a word to the Thane about the marks or the eels or the boy. The little-door burn had closed at the first edge of cold and had not opened since. The boy had been carried through two markets under this roof's rule and knew its shape now: a kindness that would not stand where an account lay open, and a silence with a wall built into it. He did not expect it to be different here. He was wrong about the difference, though he would not know it until the hand went to the cord.

A woman with her arms in a tub of salted perch tipped her head toward the basket at Sella's knee. “And that,” she said. “Whose is that, then?”

Mira did not look down. “Warm,” she said. “That's all I'll say on a slate, and I'm saying it off the slate. You want a number? There's no number in it. You want a house? Ask the fish.”

The woman laughed, said something about Fishstairs women and their secrets, and went back to the perch. That was the end of it. Mira shifted her weight so that her back was to the corner of Berr's slate, which had been left open on the stall board, and went on filleting with the knife, and the boy lay with the river wind coming into his ears and took the not-looking in beside the others he had learned to read. One man in a market had priced him without once looking into the basket, and one woman now kept him without once looking into it, and the two not-lookings had the same shape and were opposite things, and he had no word in \lore{Grathic} for either of them. He had eight days in this country and could not have asked a question in it if his life had depended on the asking, which it had, twice, without anybody knowing it.

The cold got into him through the fold. A thread of drool froze at the corner of his mouth in the river wind and stayed there, and his head went too far toward the light and made the small wet sound that was not a word, and nobody at the stall happened to be looking, so nobody put it back. His hand outside the wool went up an inch toward the bowl of eel-bones at Mira's right and came down. Two paces, then three. He had reduced a shoulder in a room the size of this stall with three people holding the light and nobody telling him what to do, and here he could not lift a bone from a bowl. The joke was there, as exact and as cheap as ever, and it cost him something and he did not laugh.

\scenebreak

Berr's voice came back from the top of the steps, given out to the stall and not to anybody in it. “Three days to \lore{the Five},” he said, and the voice had counted what it had and did not care to say more. He passed his finger over the figure on his slate the way a man might pass a finger over a thing he had set down carefully and would not now disturb. The boy had heard the number go around the district for days as a closed fist, and here it was let out into the market air in a man's ordinary voice and was still nothing he could spend. He knew the shape of the withholding. He had sat at the head of a table where a number was given out with the exactness of somebody who has decided that the room need not know what it counts.

He felt rather than saw Mira's hand go to the cord at her waist. Her thumb found a knot in it. It was a small movement, the hand going in under the apron and the thumb coming to rest and staying, and she went on talking the whole time, and the talking was about eels and the price of them and the weather on the river. The boy saw the hand go and saw the thumb settle, and he did not see the knots, and he did not know that there were knots to see. He knew only that a number had been said and a hand had gone to a place on her person, and that the place was not a slate and could not be taken from her by a man with a clean thumb. She did not say why she touched it. She did not say what it held. There was nothing in the movement that a collector would have recognised as an account, which was plainly the whole point of it.

The house he had come out of that morning kept its accounts in a roll nobody showed and a mark nobody touched except before an opening, and that keeping was a wall: a man could stand against it and be warm and never know what was written on the other side. This woman kept hers on her person, in a material no Court could confiscate, and made a small visible performance of not telling, which was a different order of thing. A wall could not be watched. A performance could. He did not trust her. He could not place her. He watched the way he had learned to watch, which was the only thing he had ever been able to do and the only thing that had never once given him a hand to act with.

He had known a hospital that kept two records: the one on the screen, which could be called up and corrected and signed by a man who had not read the note under it, and the small one a senior nurse carried in her pocket in pencil, with the things that had actually happened written small enough to hide beneath a thumb. He had never once asked to see the pocket record. He had been glad it existed and had been afraid of what it said about him, which was the same feeling from two ends. This woman's cord was that kind of record, and the boy lay under the fold with eight days of life and forty years of knowing what a second record was for, and could not have told either of them a thing. He noticed all of it. He could not have entered a single line.

He filed the day's new objects where he filed the others, side by side and joined to nothing: the eel-bones flicked and put away, the hand going under the apron, the two marks that were short and the four eels that were short, the not-looking that kept where the other not-looking had priced. He had the two ridges of \lore[the Ash]{Ash} and the charred slate and the barrow of grey. He had the death-house bench with four and four and a knot the width of a thumb, and the interval where the pot on the bench had not given its \lore{the flat note}. He knew \lore{the full note} from its flat shadow and had no way in eight days earned to make anybody produce one. The cord's movement lay beside the rest of it without a name and without a use, and he kept it carefully, the way he kept the interval of a sound that had never happened.

He had noticed the shape of the stall without deciding to notice it. A bench along the front of it, and a draught coming under the bench off the wet stones, and a space at one end of the bench where a slate should have lain and did not, because the Thane had taken his up and gone. A pot of eelweed stood where a pot stood in three other rooms he had been carried into. The arrangement followed him from house to house and then out of doors and down to the water. A bench by the door, a draught under it, a pot on the bench, and a gap in the count, and here the gap was a space on a board where an account had been set down and lifted away, leaving the shape of it in the cold. He set the shape beside the others and made nothing of the setting, because making anything of it would have been a deciding and he had no standing to decide. It was the third room in as many days that had arranged itself this way, and each time he had been carried in and set down where he could see the arrangement and not be part of it.

Mira's thumb came away from the cord and she turned back to the eels, brisk about it, her face to the stall and the fish, and the movement of the hand was the only thing in the whole market that had been done for a reason he could almost see and could not name. Sella lifted him against her hip and said something to the wool about the wind, and did not say that the woman at the stall was her sister's sister, or that the marks in the basket were short, or that the boy under the fold had been carried three doors past this place on his first morning and would be carried past it again. The boy lay at the edge of the steps with his hand outside the wool going up an inch toward the cord and coming down. Two paces, then three. The gulls went on over the ice, crying, and the river went on under its skin, and the north bridge kept its shadow on the steps.
